% 来源：OEBPS/Text/Section0001_split_006.xhtml#sigil_toc_id_40

\BookChapter{28}
\section[第28章\q 自由：创造终身收入计划]{第28章\\ 自由：创造终身收入计划}

\begin{BookQuote}
终身收入流，对退休生活的幸福很关键。

\hfill ——《时代周刊》，2012年7月30日
\end{BookQuote}

\begin{BookQuote}
我有足够的钱，可以舒舒服服地退休，享受余生。问题是，我的余生不多了，下个星期我就要死了。

\hfill ——匿名
\end{BookQuote}

1952年，埃德蒙·希拉里带领探险队，第一次成功地登上珠穆朗玛峰。大家以前一直认为这是不可能完成的任务。英格兰女王马上封他为爵士，以表彰他的惊人成就，从此他变成了埃德蒙·希拉里爵士。

尽管他取得这么大的成就，很多人还是认为，埃德蒙·希拉里爵士也许并不是第一个到达珠穆朗玛峰顶峰的人。事实上，大家普遍认为，乔治·马洛里也许是第一个登顶珠穆朗玛峰的人，比希拉里要早上30年。

那么，如果乔治·马洛里1924年就到达珠穆朗玛的顶峰，为什么是30年之后才登顶的埃德蒙·希拉里得到了所有的荣誉——包括被女王册封为爵士？

因为埃德蒙·希拉里不仅到达了顶峰，他还成功地从顶峰下来了。乔治·马洛里就没有那么幸运了。就像绝大多数死在珠穆朗玛峰上的人一样，结果证明下山才是致命危险的。俗话说得好：上山容易下山难。

\subsection{你投资为了什么？}

我经常问别人：“你投资为了什么？”

答案各种各样，什么样子的都有：

“收益。”

“增长。”

“资产。”

“自由。”

“乐趣。”

我很少听到那个最重要的答案：收入！

我们都需要有一份收入来维持生活。有一个持续不断的现金流，每个月都出现在我们的账户上，像时钟转动一样可靠。你能想象你再也不用担心怎么支付每个月的水费、电费、天然气费、手机费的情景吗？你能想象你再也不用担心你的钱会花光了吗？你自由自在、快快乐乐地到处旅行，根本不用担心钱的事。每个月打开股票账户对账单的时候，你不用再祈祷市场挺住，内心非常安宁平静，慷慨大方地承诺捐钱给你的教堂或者你最喜欢的慈善团体，心里不用担心以后每年捐的钱从哪里来。我们凭直觉都知道：收入就是自由。

就像电影《勇敢的心》中的一幕，梅尔·吉布森在山顶上大喊：“收入就是自由。”

缺少收入，你心里就有压力。缺少收入，你的日子就很难过。缺少收入，对于你和你的家庭来说，是不可能接受的结果。

杰弗里·布朗博士，在最近一期《福布斯》杂志的文章中，对这一点说得很明白：“收入就是结果，没有足够的收入，你的退休生活就不安全。”

富人都知道，他们那些资产（股票、债券、黄金等）的市场价值总是在上下波动。但是你不能拿那些资产“当钱花”，你能消费的只有现金。2008年就出现了这样的事情，很多人拥有的资产（特别是房地产）市场价格大跌，他们很想卖出，却怎么也卖不掉。他们在资产上很“富有”，在现金上却很“贫穷”。因为资产不能变现，到期的债务还不了，只能破产。你始终要牢牢记住：收入就是结果。

在这一部分的末尾，你会有具体的工具来锁定你期望得到的具体收入水平。我称之为“收入保险”。这是一种有保障的方式，非常确定地知道，你会有一张终身支票，将来再也不用为钱而工作——你绝对可以确定你永远花不完你的钱。猜猜你需要做什么？你只需要决定你想让你的终身收入支票什么时候开始兑现。

要得到这样的收入保险，有很多办法。我们来看几个对于你来说比较合适的办法，它们让你获得收入保险。

有一个锁定收入的结构还有很多其他的好处。这是地球上唯一能给你带来以下好处的金融工具：

\begin{itemize}
\item 你的本金百分之百地保证安全。\footnote{\label{fn:pdf-28-1}保险保障协会提供保护给保单持有人和保单受益人，以避免签发保单的保险公司破产而无法兑现自己的保险合同责任，美国各州、哥伦比亚特区和波多黎各，都有保险保障协会。保险公司按照法律要求，在其得到执照允许发展业务的州，必须成为所在州保险保障协会的成员。每个州都有自己规定的最高保护金额限制，大多数州最高保护金为30万\textasciitilde{}50万美元。}（你不可能亏钱，你可以完全控制你交的钱。）
\item 只有上行收益，没有下行风险：你的账户价值增长会和市场涨幅相关联。如果股市上涨了，你可以参与分享一些收益。但即使市场下跌了，你也不会损失一分钱。
\item 你的投资增值可以递延纳税。（还记得那个价值翻倍的案例吗？税收效率决定你最终是拥有28466美元，还是拥有超过100万美元。）
\item 一份有保障的终身收入流：你有控制权，可以决定什么时候开始领取这份收入。
\item 记住这一点：付给你的收入，如果结构安排合理的话，是可以免税的。
\item 没有年度管理费。
\end{itemize}

你能得到所有这些好处，只要运用一个有2000年历史的金融工具的现代版就行了！这怎么可能？我可以肯定，这听起来简直太好了，不可能是真的，但是相信我，这确实是真的！我就在用这种投资方式，使用它之后让我非常激动，我很高兴能给你详细地讲述如何操作。

本书一直在强调，你盼望达到的财务未来非常像是你在攀登珠穆朗玛峰。你要工作几十年，不断地积累财富，直到达到一个关键的规模（爬到顶峰），但是这只是故事的前一半。如果财富积累到关键规模，却没有一个计划和策略，如何把这些资产转化为收入，在你的余生一直持续不断地维持你正常的生活？在这样的情况下，如果收入断了，你就会像乔治·马洛里一样死在下山路上。

\subsection{新时代}

毫无疑问，我们现在处在一个完全没有人探索过的海域里，这个时代的情况绝对可以说前所未有。在过去30年里，退休这个概念发生了彻底的改变。直到20世纪80年代末，在美国，62\%以上的工作者都有退休金计划。你还记得吗？你上班的最后一天得到一块金表，还有一份有保障的终身收入支票，你一退休就能领到第一张。现在，除非你是给政府工作，退休金已经很少有了，它成为历史遗迹，像个金融界的恐龙。现在不管是好是坏，你都是你自己的船的船长。你的钱能不能持续维持生活，最终是由你自己负责的。这是你自己要承担的一个重大责任，也是一个重大负担。自己管理自己的养老金投资，意味着你下海了，下到了股海之中，股市的波动很多时候比海浪还要大，加上基金过高的收费，通货膨胀，身体出现“意外”花去大笔医疗费用，你很快开始明白，为什么会有这么多人面临巨大的退休危机。很多人，包括你的邻居和同事，都要面对很可能出现的问题：人还活着，钱已经没了。特别是现在，人的平均寿命比以前要长得多。

\subsection{现在的80岁相当于过去的50岁}

时间长又丰富多彩的退休生活，这个概念只是最近几代人才有的。我们前面讨论过退休制度方面的内容，1935年富兰克林·罗斯福总统创立了社会保险制度，那时候人的平均寿命预期只有62岁。但是年满65周岁才能领到社会保险，所以在社会保险制度刚开始实施那些年，其实只有少数活得比一般人更加长的人，才能真正享受到社会保险的福利。

那个时候的社会保障体系在财务上很稳定，因为平均40个交纳社会保险的工作者来养活一个领社会保险的退休者。这就像是有40个人拉车，只有一个人坐车。到2010年，这个比例下降为每1个坐车的人，只有2.29个人拉车。什么时候开始这件事难住了美国政府？

现在美国男性的平均寿命预期为79岁，女性为81岁，一对已婚夫妇，其中至少有一个人有25\%的机会能活到97岁。

\subsection{但是等等，以后的人还会活得更久！}

你可能会活得比平均寿命预期还要久。想一想最近30年来我们的科技发展进步有多么大。从软盘到纳米科技。现在的科学家正在用3D打印制造新器官，就像魔术师变戏法一样。研究者能用人的细胞轻轻地在你的皮肤上刮一下，就可以用这些细胞作为基础“打印”出一个新耳朵、新膀胱、新气管！\footnote{\label{fn:pdf-28-2}安东尼–阿塔拉，威克弗里斯特再生医学研究所的主任，过去十几年来，一直在创造和植入这种器官。}原来科幻小说里面幻想的东西，现在都成真了。后面我们会直接聆听我的朋友雷·库兹韦尔所言，他是我们这个时代的爱迪生，最近他担任了谷歌工程部门的负责人。我问他：“生命科技的进步会如何影响人类寿命预期？”他说：

21世纪20年代，人类会找到方法改变人类的基因。不仅“人工设计婴儿”会成为现实，而且会出现人工设计婴儿潮，因为所有人的人体组织和器官都可以再生，可以把一个人的皮肤细胞转化成每一种其他类型细胞的年轻版。人们可以“重新编程”自己的生物化学机制，远离疾病和衰老，大幅度地延长寿命预期。

对于美国婴儿潮世代来说，这些话听起来太让人激动了！皱纹走开吧！我们马上都可以喝到青春不老泉了。

但是，对于我们的退休生活来说，钱的问题很明显就大多了。我们原来以为只能活到80岁，也是按照这个计划攒的退休养老金。现在科技进步让我们的寿命变长，相比之下退休要花的钱也多得多了，原来攒的钱需要维持的退休生活年数也更长。你想象一下，如果雷·库兹韦尔说的是对的，美国婴儿潮一代出生的人，现在60\textasciitilde{}70岁，将来会活到110岁，甚至120岁，那么他们还有40\textasciitilde{}60年的退休生活呢！再想象一下，这样的生命科技进步会改变21世纪初出生的那一代人的寿命预期。如果你将来还要活110\textasciitilde{}115年，你需要攒多少钱才够养老呢？没有什么比一份有保障的终身收入更重要了。只要你活着，就有一份终身收入，绝对不会人活着钱却没了，这对于你来说，会是最好的资产。

\begin{BookQuote}
年轻的时候，我认为钱是人生最重要的东西。现在，我老了，我承认事实就是这样的。

\hfill ——奥斯卡·王尔德
\end{BookQuote}

\subsection{4\%法则完了}

20世纪90年代初期，美国加利福尼亚州有个投资理财规划师，提出了“4\%法则”。这个4\%法则的要点是，如果你想让你的钱能够持续一生花不完，你可以每年取出来4\%，前提是你要建立一个“平衡投资组合”，60\%投资股票，40\%投资债券。你可以每年增加取现的比例，以抵消通货膨胀。

“确实，在过去相当长的时间内，4\%法则表现很好。”2003年《华尔街日报》上有一篇文章回顾道，文章标题是《跟4\%法则说再见》。为什么4\%法则会突然被抛弃？因为4\%法则是在1990年左右被提出来的，当时美国国债支付的年利率超过4\%，股市正好是一个大牛市！如果你是在2000年1月退休的，按照传统的4\%法则来操作，那么到2010年你的资金就会亏损33\%，100万美元只剩下67万美元，按照美国十大基金公司之一普信集团的研究，你的钱按照4\%法则操作，只有29\%的概率能够支撑你余下的寿命。我换个更加直截了当的说法是，你有71\%的概率是人活着而钱没了。我们大多数人最不愿意同时体验的两件事是：人老了，钱没了。

现在，我们生活在一个全球各国都在压低利率的世界中，事实上这就是一场针对储蓄者的战争，可以肯定的是，这就是一场针对老年人的战争。利率接近于零，谁还能安全地过退休生活呢？退休的老年人必须冒险进入不安全的投资领域，给自己多年攒的钱找到适当的回报，就像那些渴得要死的羚羊，必须冒险去鳄鱼密集的水域喝水，没办法，不喝水它们肯定只能渴死。股市危机四伏，但是那些老年人也没办法，他们需要有正向回报来维持生活，每个月的水费、电费账单必须要付，但投资股市会让他们更容易受到伤害。

\subsection{群聚效应被摧毁}

无论什么人告诉你或者兜售给你投资方法，没有任何一个基金经理、经纪人或投资理财顾问能够控制我们的钱足够维持一生。这就是金融世界肮脏的小秘密，只有少数专业人士知道其中的决定因素。而在少数知道这个秘诀的人里，只有极少数人愿意公开地说出来。我说话一向直率，那天我与投资界传奇人物约翰·博格坐到一起交谈，我直接把这个问题抛了出来。

还记得约翰·博格吗？他是世界上最大的指数基金公司先锋公司的创始人，他说话很直接。那天我跑到宾夕法尼亚州，在他办公室里谈了4个小时，我抛出来那个肮脏的小秘密，他也毫不掩饰他的观点和想法。“有些东西，我不愿意说，但是所有这些东西有时候就像赌博一样，完全靠运气：你什么时候出生，你什么时候退休，你的孩子什么时候上大学，你根本没办法控制。

他说的赌博是怎么一回事？

其实就是说有很大的运气成分在里面：关键要看你退休的时候，市场表现怎么样。如果你在1995年前后退休，你就是一个快乐幸福的度假“露营者”。如果你在2005年左右退休，你就是一个无家可归的“露营者”。博格自己亲口说过，那是在2013年年初的一次电视节目上接受采访，下个10年我们要做好准备，股市可能会有两次大跌，最大跌幅会有50\%。我的天哪！但是对他这样的预测，我们当时也许不应该感到吃惊。2000—2009年，美国股市经历过两次下跌接近50\%。不要忘记，如果你亏了50\%，那么你需要再赚100\%才能赚回你原来的本金。

这个肮脏的小秘密就是，我们都面临的风险是收益连续性的毁灭性概念。听上去很复杂，但其实很容易懂。从本质上讲，你退休后最开始几年的投资业绩的好坏将会决定你后面好几十年投资业绩的好坏。简单地说，退休前面几年决定后面的几十年。如果你退休之后前几年遭受了投资亏损，这完全是运气不好，但是你以后把钱赚回来弥补亏损，并且赚到钱让你退休生活衣食无忧，这样的概率低到像跌落悬崖一样，可以说永远不可能发生。

你可以做到把每一件事都做对：找到一个受托人投资理财顾问，减少你的费用，投资很有税收效率，建立一个财务自由基金。

但是当你爬到顶峰，实现财务自由之后，你要下山就是从赚钱变成花钱，要从你的投资组合中获得收入以维持生活了。你如果在退休后前几年里碰到一年的投资业绩非常糟糕，就会很容易打乱你的整个计划，使生活变得一团糟。若是更加不幸，碰上几年投资业绩都比较糟糕，你可能会发现，自己都退休好几年了，又不得不回去重新工作，卖掉你买好的度假屋。这听起来是不是变化太大了？让我们看一个假设的案例，看一看收益先后顺序所带来的风险，长期而言对我们的影响有多大。

\subsection{约翰咬了狗}

约翰咬了狗，狗咬了约翰，都是同样5个字，但是，先后顺序安排不同，就有了完全不同的意思，特别是约翰这两个字，放置的顺序不一样，其意思差别会很大。

约翰现在65岁，攒了50万美元的养老金（远远超过美国人的平均水平），他准备要退休了。跟大多数即将退休的美国人一样，约翰也正在持有一个“平衡的”投资组合（60\%股票加上40\%债券）。我们前面从瑞·达利欧那里知道了，这样的风险配置其实并不平衡，因为利率都这么低了，接近于零，4\%法则根本不适用。约翰决定，他每年从50万美元的养老金里取出来5\%，就是2.5万美元，来满足日常生活的基本需要，再加上他能够领到社会保险，这些钱他应该够花了，可以生活得不错。他还必须每年提高取现比例（每年提高3\%），才能赶得上通货膨胀，因为物价一上涨，每一年同样数量的钱能买到的东西就减少了。

约翰的运气不好，退休前几年遇到几次股市下跌，股票投资亏了不少。事实上，他称自己退休之后的生活为黄金年代，结果前三年股市就连续大跌三年。这可不是什么美好生活该有的开始。

退休之后才过了短短5年，约翰原来攒的50万美元的退休金就亏掉了一半。股市下跌时，他还要取现，只能亏本卖出股票换取现金。没办法，每个月好多账单要付，这真是雪上加霜。由于你的投资金额减少了，如果股市反弹，你赚的钱也就减少了，回本就更难了。但没办法，生活还得继续，账单必须现在就付。

% 来源：OEBPS/Text/Section0001_split_006.xhtml → image_72e490c7cee64829b27f5bca9b708466.jpg
\begin{longtable}{@{}C{\dimexpr 0.10000\linewidth-1.50000\tabcolsep\relax}C{\dimexpr 0.26000\linewidth-1.50000\tabcolsep\relax}C{\dimexpr 0.30000\linewidth-1.50000\tabcolsep\relax}C{\dimexpr 0.34000\linewidth-1.50000\tabcolsep\relax}@{}}
\caption{约翰的投资}\label{tab:28-1}\\
\toprule\noalign{\vskip 3pt}
{\heiti 年龄} & {\heiti 假设股市盈利或亏损（\%）} & {\heiti 每年年初取现（美元）} & {\heiti 每年年初养老金（美元）} \\
\midrule
\endfirsthead
\multicolumn{4}{r}{续表}\\[0.6em]
\toprule\noalign{\vskip 3pt}
{\heiti 年龄} & {\heiti 假设股市盈利或亏损（\%）} & {\heiti 每年年初取现（美元）} & {\heiti 每年年初养老金（美元）} \\
\midrule
\endhead
\bottomrule
\endfoot
\bottomrule
\multicolumn{4}{@{}p{\linewidth}@{}}{\footnotesize 平均收益：8.03\%。总取款额：580 963美元。}\\
\endlastfoot
64 &  &  & 500\,000 \\
65 & −10.14 & 25\,000 & 500\,000 \\
66 & −13.04 & 25\,750 & 426\,839 \\
67 & −23.37 & 26\,523 & 348\,766 \\
68 & 14.62 & 27\,318 & 246\,956 \\
69 & 2.03 & 28\,318 & 251\,750 \\
70 & 12.40 & 28\,982 & 228\,146 \\
71 & 27.25 & 29\,851 & 223\,862 \\
72 & −6.56 & 30\,747 & 246\,879 \\
73 & 26.31 & 31\,669 & 201\,956 \\
74 & 4.46 & 32\,619 & 215\,084 \\
75 & 7.06 & 33\,598 & 190\,084 \\
76 & −1.54 & 34\,606 & 168\,090 \\
77 & 34.11 & 35\,644 & 131\,429 \\
78 & 20.26 & 36\,713 & 128\,458 \\
79 & 31.01 & 37\,815 & 110\,335 \\
80 & 26.67 & 38\,949 & 95\,008 \\
81 & 19.53 & 40\,118 & 71\,009 \\
82 & 26.38 & 36\,923 & 36\,923 \\
83 & −38.49 & 0 & 0 \\
84 & 3.00 &  &  \\
85 & 13.62 &  &  \\
86 & 3.53 &  &  \\
87 & 26.38 &  &  \\
88 & 23.45 &  &  \\
89 & 12.78 &  &  \\
\end{longtable}


从70岁开始，约翰在股市上享受稳定的上涨年份，但是损害已经造成。要弥补亏损的道路实在太崎岖了。你知道的，亏损一半之后，要再赚到翻倍的收益才能回本。等到他76岁左右，他已经看得很明白了，他的钱再过几年就要花完了。等到他83岁时，不幸又碰到一次股市大跌，他的股票账户卖光了。最终，总的来看，他原来攒了50万美元的退休金账户，退休生活18年只从中累计取现580930美元。换句话说，退休后他又继续投资了18年，可是总共只赚了8万美元。

但是其中最疯狂的事情是：在约翰跌跌撞撞一路下山（退休生活）的时候，股票市场平均每年上涨超过8\%。这是相当好的业绩，不管是用谁的标准来看都相当好。

问题就在这里，市场并不是每年都能给你相同的平均收益率，他有的年份收益率比平均水平高，有的年份收益率比平均水平低，只是长期下来这些年份整体来看是这样的平均收益率而已。（记住我们在前面讨论过真实收益和平均收益之间的差别，就在第二部分第4章，“希望”在有些年份里你不会遭受损失，因为你承受不起，这并不是一个有效的策略，无法确保你的财务未来安全无忧。）

\subsection{先后顺序调过来}

苏珊女士和约翰的情况几乎完全相同，她也是65岁，攒了50万美元。她和约翰一样，每年要提现5\%，就是2.5万美元作为她每年的收入来维持退休之后的生活，她需要每年略微增加一点投资变现的比例，这样调整是为了应对通货膨胀带来的物价上涨。为了真切地描述这个概念，我们会对两个人用相同的投资收益率，我们只是把这些收益率的顺序颠倒一下，把第一年和最后一年颠倒顺序，第二年跟倒数第二年颠倒顺序，以此类推。

只是颠倒收益的顺序，就会让苏珊的退休生活体验完全不同于约翰。事实上，等到89岁的时候，苏珊累计投资变现超过90万美元，这些都是每年变化的收益作为当年收入用来消费的，但是花了这么多钱，她的投资账户上还剩下1677975美元！苏珊在这个世界上生活的时候，从来没有担忧过钱的事。

% 来源：OEBPS/Text/Section0001_split_006.xhtml → image_df105b1f265b49f8a188f81dd9928fbf.jpg
\begin{longtable}{@{}C{\dimexpr 0.10000\linewidth-1.50000\tabcolsep\relax}C{\dimexpr 0.26000\linewidth-1.50000\tabcolsep\relax}C{\dimexpr 0.30000\linewidth-1.50000\tabcolsep\relax}C{\dimexpr 0.34000\linewidth-1.50000\tabcolsep\relax}@{}}
\caption{苏珊的投资}\label{tab:28-2}\\
\toprule\noalign{\vskip 3pt}
{\heiti 年龄} & {\heiti 假设股市盈利或亏损（\%）} & {\heiti 每年年初取现（美元）} & {\heiti 每年年初养老金（美元）} \\
\midrule
\endfirsthead
\multicolumn{4}{r}{续表}\\[0.6em]
\toprule\noalign{\vskip 3pt}
{\heiti 年龄} & {\heiti 假设股市盈利或亏损（\%）} & {\heiti 每年年初取现（美元）} & {\heiti 每年年初养老金（美元）} \\
\midrule
\endhead
\bottomrule
\endfoot
\bottomrule
\multicolumn{4}{@{}p{\linewidth}@{}}{\footnotesize 平均收益：8.03\%。总取款额：911 482美元。}\\
\endlastfoot
64 &  &  & 500\,000 \\
65 & 12.78 & 25\,000 & 500\,000 \\
66 & 23.45 & 25\,750 & 535\,716 \\
67 & 26.38 & 26\,523 & 629\,575 \\
68 & 3.53 & 27\,318 & 762\,140 \\
69 & 13.62 & 28\,318 & 760\,755 \\
70 & 3.00 & 28\,982 & 832\,396 \\
71 & −38.49 & 29\,851 & 827\,524 \\
72 & 26.38 & 30\,747 & 490\,684 \\
73 & 19.53 & 31\,669 & 581\,270 \\
74 & 26.67 & 32\,619 & 656\,916 \\
75 & 31.01 & 33\,598 & 790\,788 \\
76 & 20.26 & 34\,606 & 991\,981 \\
77 & 34.11 & 35\,644 & 1\,151\,375 \\
78 & −1.54 & 36\,713 & 1\,496\,314 \\
79 & 7.06 & 37\,815 & 1\,437\,133 \\
80 & 4.46 & 38\,949 & 1\,498\,042 \\
81 & 26.31 & 40\,118 & 1\,524\,231 \\
82 & −6.56 & 41\,321 & 1\,874\,535 \\
83 & 27.25 & 42\,561 & 1\,712\,970 \\
84 & 12.40 & 48\,383 & 2\,125\,604 \\
85 & 2.03 & 45\,153 & 2\,339\,923 \\
86 & 14.62 & 46\,507 & 2\,341\,297 \\
87 & −23.37 & 47\,903 & 2\,630\,297 \\
88 & −13.04 & 49\,340 & 1\,978\,993 \\
89 & −10.14 & 50\,820 & 1\,677\,975 \\
\end{longtable}


两个人退休之时攒了同样多的退休金，同样的投资取现策略，退休生活过得差别却很大，简直是天壤之别：一个人老了变成了穷人，另外一个人则实现了绝对财务自由。

更让人感到迷惑不解的是：这25年间两个人的平均投资业绩也相同，都是每年8.03\%。

这怎么可能？因为平均业绩就是总的收益除以年数。

没有人能够预测，下一个拐角会发生什么。没有人知道市场什么时候会上涨，什么时候会下跌。

现在，想象一下，要是约翰和苏珊都有收入保险。约翰会避免得胃溃疡，因为他知道，随着他的账户缩水，他有一个有保障的收入支票就在彩虹那头等着他。苏珊会有更多的钱可以花，她想怎么花就怎么花。也许可以再多买一个度假屋，给她的孙子孙女更多钱花，捐赠更多钱给她最喜欢的慈善事业。收入保险的价值无论怎样都不会被高估。搭配全天候投资组合，你就有了最强有力的双保险。

\subsection{一个人拿了6个学位}

你也许还记得，在本书的开头，我跟你提过一个沃顿商学院的教授戴维·巴贝尔博士，我认识的成千上万个人里面，他不仅是受教育水平最高的人，也是一个灵魂上温柔体贴但信念上坚如钢铁的人。戴维·巴贝尔喜欢别人称呼他时在名字前加上博士或教授。

我们来快速回顾一下戴维·巴贝尔的成就。他有6个学位：经济学学士、国际金融专业工商管理硕士、金融博士、辅修食品和资源经济学博士、热带农业博士、拉丁美洲研究博士。戴维·巴贝尔在加州大学伯克利分校和沃顿商学院教授投资课程超过30年。他是高盛集团养老金和保险研究分部主任。他曾经为世界银行工作，提供咨询的客户包括美国财政部、美联储、美国劳工部。戴维·巴贝尔对养老金问题的研究，可以说跟篮球之神乔丹对篮球的了解一样。

戴维·巴贝尔写过一篇研究报告，得到的评价褒贬不一，两极分化。在这篇报告里，他分析的是自己的个人退休养老金计划。戴维·巴贝尔即将退休，他希望能够找到一种策略，给他带来内心的宁静，也带来一份有保障的终身收入。他始终牢记，收入就是结果。他还明智地考虑了其他因素，例如不愿意自己年纪很大了还要做出复杂的投资决策，他考虑了自己能找到的所有选择，也调用他对风险和市场的巨大知识库，他甚至咨询了他的朋友和以前他在华尔街工作时的同事，对各种策略进行比较分析。最终戴维·巴贝尔决定，他辛辛苦苦赚来的退休养老金，最好的投资去处就是年金。

啊！等等。

戴维·巴贝尔的华尔街朋友称年金为“年杀”，他这个资深投资专家怎么能把自己的养老金托付给年金呢？“年杀”这个词是由那些证券公司的经纪人炮制出来的，用来形容他们的有些客户把钱抽离股市，找个年头非常长的保险公司，购买年金以保证得到一份稳定的终身收入。证券公司经纪人觉得，这是一个不可逆转的决策，因为你把股票账户上的钱取走买保险公司的年金，这些钱就永远不会回到股票账户上了，他们就再也不能从你的股票投资上赚到佣金收入了。其实死掉的是证券公司的盈利。

好好想一想，最近一次经纪人跟你讲如何创造一份终身收入是在什么时候？很可能从来没有过。证券公司一般没有兴趣推广把钱从股市拿出来的相关概念，对于证券公司来说，一听说你要退出市场，它们就着急了。令人感到嘲讽的是，只要你不销户离开股市，你对于证券公司来说就代表着一份终身收入。

\begin{BookQuote}
美国人至少应该把他们一半的退休储蓄转成年金。

\hfill ——美国财政部
\end{BookQuote}

杰弗瑞·布朗博士对创造终身收入计划，肯定认识得更深刻。他是美国财政部和世界银行的顾问，中国政府邀请来评估未来社会保障战略的专家之一。他还是时任美国总统指定的美国社会保障顾问委员会7位成员之一。

杰弗瑞·布朗博士职业生涯的大部分时间，都花在了研究如何给人们提供一份终身收入上。那么他是怎么解决这个问题的？那就是利用年金这个我们所拥有的最重要的投资工具之一。

我和杰弗瑞·布朗交流了三个小时，我们谈得非常愉快，我们谈的主题就是收入计划，让他困惑不解的是，大多数投资理财规划师在和客户交谈时经常会省略收入。这怎么可能？收入保险只在大多数投资理财规划师的办公室里，在同事之间才会讨论，却没有投资理财规划师和客户讨论，也没有包含在401(k)计划的投资选项中，而401(k)计划是美国人最主要的退休金投资工具。

我问杰弗瑞·布朗：“普通人应该如何保护自己，让自己真的拥有一份终身收入，即使自己活的比以前的人长得多，也不用担心人还在，而钱没了。美国现在都是年满65岁退休，按照寿命预期，退休后还要活上二三十年，自然需要二三十年的收入来支撑退休后的生活，但是他们的投资理财规划却没有持续那么长。寿命有那么长，钱却不够花那么长。你有什么解决方法吗？”

“托尼，好消息是我们确实知道如何来解决这个问题。”他说，“我们只需要让大家改变为退休养老攒钱的思考方式就行了。在‘经济学家的土地’上，有一种产品就放在那里，我们称之为年金。基本你都可以到保险公司那里去说，‘你知道吗？我要把我的钱带过来，放在你那里，由你来管理，让这些钱不断升值，作为回报，你要每个月付给我一份收入，只要我还活着就有收入’。有个很容易的方法让你一下子就能理解年金这个产品，其实年金做的事和社会保险一模一样。大家都知道，上班工作时，每个月你都要交社会保险，之后等到你退休了，你就能每个月从社会保险那里领到一份收入，只要你活着就能一直领下去。但是光靠社会保险那些收入并不足以维持你想要的退休生活。买了年金，相当于你给自己多存了几份社会保险，靠自己来扩大终身收入。”

杰弗瑞·布朗和他的团队做了一项研究，比较用不同的方式来描述或者“框定”年金，塑造年金的形象，是如何完全改变人们对年金的需要或者渴望的。

第一，他们像证券公司的股票经纪人那样描述年金：年金作为一个“储蓄”账户，或者只有相对于股市而言收益率较低的投资。一点儿也不奇怪，只有20\%的人觉得年金有吸引力。这听起来是不是很熟悉？你总能听到证券公司的经纪人说“年金是个糟糕的投资”。

但是，后来他们只改了几个词描述年金能够实际带来的真正好处，年金给人的形象马上就改变了。他们这样描述：年金作为一种工具能给你带来有保障的收入，而且一直持续，你的寿命有多久，这份收入就有多久，结果超过70\%的人听到之后都觉得年金很有吸引力。谁不想有一份收入保险呢？万一你花光了积蓄，不要慌，你可以马上开始领取收入，这可是救命的钱。也许你的生活开支会比你预想的大；也许你发生意外要做紧急手术，要花一大笔钱；也许股市的走势和你预期的不一样，跟你的支出需求不合拍——你要花钱的时候，股市反而在下跌。这个时候，如果你知道自己有一份终身收入支票，只要打个电话就能开始每个月领取收入了，那么它会是一份多么及时的大礼啊。

现在的金融行业有很多革命性的创新，已经创造出一套全新的年金投资机会。其中很多投资可以支付给你一定比例同期股市的上行收益率，却不用承担任何股市的下行损失。现在的年金再也不是原来你爷爷买的那种老式年金了。翻到下页，我来给你展示5种类型的年金，它们可能会改变你的人生。
